\documentclass[11pt,a4paper]{article}
\usepackage[utf8]{inputenc}
\usepackage{amsmath,amssymb,amsfonts,amsthm}
\usepackage{graphicx}
\usepackage{hyperref}
\usepackage{booktabs}
\usepackage{geometry}
\geometry{margin=1in}

\newtheorem{theorem}{Theorem}[section]
\newtheorem{definition}[theorem]{Definition}
\newtheorem{proposition}[theorem]{Proposition}
\newtheorem{lemma}[theorem]{Lemma}

\title{\textbf{Theoretical Formalization of Automatic Prompt Optimization (APO) in Multimodal Environments: A Graph-Based Textual Interface Approach}}
\author{AI Research Assistant \& Collaborator}
\date{\today}

\begin{document}

\maketitle

\begin{abstract}
Automatic Prompt Optimization (APO) has emerged as a crucial paradigm for maximizing the performance of Large Language Models (LLMs) and multimodal reasoning systems without weight updates. However, formalizing APO in multimodal environments---where inputs span visual, auditory, and textual domains---remains challenging due to the complex discrete search spaces and heterogeneous textual interfaces. In this paper, we propose a rigorous theoretical framework modeling APO as a discrete optimization problem over graph-based textual interfaces. We define the multimodal interface graph, formalize the prompt optimization objective using expected utility over task distributions, and analyze the convergence properties of iterative discrete optimizers. Furthermore, we outline experimental validation pathways and discuss implications for robust, interpretable multimodal agent architectures.
\end{abstract}

\section{Introduction}
Large Language Models (LLMs) and multimodal foundational models (OmniLLMs) rely heavily on prompting strategies to steer their behavior across diverse tasks, such as audio-visual question answering, scientific simulation design, and code generation. While manual prompt engineering is intuitive, it is brittle, highly sensitive to minor phrasing variations, and fails to scale. Consequently, \emph{Automatic Prompt Optimization (APO)} has been developed to iteratively refine prompt tokens using feedback gradients or LLM-based critics.

Recent literature highlights that in multimodal domains (e.g., audio-visual reasoning), direct end-to-end optimization is complicated by the disparate native modalities of sensory inputs. As demonstrated by Wang et al.\ (2026), systems must rely on structured textual interfaces---such as direct inference models, captioner-reasoner pipelines, or multi-agent orchestration graphs. 

Despite empirical advances, a principled mathematical framework for APO over graph-based textual interfaces is lacking. This paper addresses this gap by formalizing:
\begin{enumerate}
    \item A graph-based representation of multimodal textual interfaces.
    \item A discrete optimization objective function mapping text-prompt spaces to expected multimodal utility.
    \item Convergence and complexity bounds for iterative prompt search algorithms.
\end{enumerate}

\section{Preliminaries and Graph-Based Textual Interfaces}
Let $\mathcal{M}$ be a multimodal foundational model capable of processing sensory inputs $x \in \mathcal{X}$ (where $x = (x_{\text{vid}}, x_{\text{aud}}, x_{\text{text}}))$ and producing an output $y \in \mathcal{Y}$. 

\begin{definition}[Multimodal Textual Interface Graph]
A Multimodal Textual Interface Graph (MTIG) is defined as a directed acyclic graph $G = (V, E)$, where:
\begin{itemize}
    \item $V = V_{\text{data}} \cup V_{\text{trans}} \cup V_{\text{reason}}$ represents the set of nodes corresponding to input data sources, feature transformers (e.g., captioners, automatic speech recognition modules), and reasoning agents, respectively.
    \item $E \subseteq V \times V$ represents the directed flow of textual and tokenized information between modules.
    \item Each node $v \in V$ is parameterized by a local system prompt $p_v \in \Sigma^*$, where $\Sigma$ is the vocabulary set.
\end{itemize}
\end{definition}

Let $\mathbf{P} = (p_{v_1}, p_{v_2}, \dots, p_{v_k})$ denote the joint prompt vector across all graph nodes $V$. The global behavior of the system under input $x$ is governed by the composite operator:
\begin{equation}
    f_{\mathbf{P}}(x) : \mathcal{X} \to \mathcal{Y}.
\end{equation}

\section{Problem Formulation: Discrete Optimization of Prompts}
Given a task dataset $\mathcal{D} = \{(x_i, y_i\}_{i=1}^N$ drawn from an unknown task distribution $\mathcal{P}_{XY}$, and a loss function $\ell : \mathcal{Y} \times \mathcal{Y} \to \mathbb{R}$, our goal is to find the optimal prompt vector $\mathbf{P}^*$ that minimizes the expected loss over the graph MTIG:

\begin{equation}
    \mathbf{P}^* = \arg\min_{\mathbf{P} \in \boldsymbol{\Omega}} \mathbb{E}_{(x,y) \sim \mathcal{P}_{XY}} \left[ \ell(f_{\mathbf{P}}(x), y) \right],
\end{equation}

where $\boldsymbol{\Omega} = \Sigma_1^* \times \Sigma_2^* \times \dots \times \Sigma_k^*$ is the discrete combinatorial prompt space. 

Because $\boldsymbol{\Omega}$ is discrete and high-dimensional, gradient-based backpropagation through model weights is infeasible. Instead, APO relies on derivative-free optimization (DFO) or LLM-guided iterative refinement (e.g., using a meta-optimizer model $\mathcal{M}_{\text{opt}}$).

\section{Theoretical Analysis of Iterative APO Convergence}
We model iterative APO as a Markov Chain over the discrete state space $\boldsymbol{\Omega}$. At iteration $t$, the optimizer proposes a candidate prompt update $\mathbf{P}^{(t+1)} \sim q(\cdot \mid \mathbf{P}^{(t)}, \mathcal{E}^{(t)})$, where $\mathcal{E}^{(t)}$ is the execution error feedback on a mini-batch of validation data.

\begin{lemma}[Monotone Improvement under Perfect Feedback]
Let $J(\mathbf{P}) = \mathbb{E}[\ell(f_{\mathbf{P}}(x), y)]$ be the true expected loss. If the prompt proposal mechanism $q$ incorporates an exact gradient surrogate via LLM feedback such that proposed prompts satisfy $\Pr(J(\mathbf{P}^{(t+1)}) < J(\mathbf{P}^{(t)})) \ge 1 - \epsilon$, then the expected loss sequence $\{J(\mathbf{P}^{(t)})\}_{t=0}^{\infty}$ converges almost surely to a local minimum in $\boldsymbol{\Omega}$.
\end{lemma}

\begin{proof}
Since $\boldsymbol{\Omega}$ is finite (bounded by maximum token length $L$ and vocabulary size $|\Sigma|$), and the loss function $J(\mathbf{P})$ is bounded below by $0$, any sequence of monotonically decreasing values on a finite state space must converge. By standard martingale convergence theorems applied to discrete search spaces with drift bounded by $\epsilon$, the process converges to a stationary distribution concentrated on local optima.
\end{proof}

\section{Discussion and Future Research Directions}
The graph-based formulation established above provides several distinct advantages:
1. \textbf{Modularity:} Errors can be isolated to specific nodes in $G$ (e.g., faulty audio transcription vs.\ flawed reasoning prompt), allowing localized prompt tuning.
2. \textbf{Scalability:} Decomposing monolithic prompts into graph-distributed prompts reduces the combinatorial search space from $|\Sigma|^{kL}$ to $\sum_{v} |\Sigma|^{L_v}$.

Future work will implement this framework empirically across multimodal benchmarks and analyze sample complexity bounds under noisy LLM-as-a-judge reward signals.

\begin{thebibliography}{9}
\bibitem{wang2026designing}
Y. Wang, T. Wang, T. Lei, D. Wang, and P. Bell.
\newblock Designing Textual Interfaces To Enhance Llm Reasoning On Audio-Visual Question Answering.
\newblock \emph{arXiv preprint arXiv:2610.20085}, 2026.

\bibitem{chen2026speccloak}
W. Chen, J. Wang, G. Wu, J. Liang, and Y. Cai.
\newblock SpecCloak: Probing Harmful Code Generation through Engineering Specifications.
\newblock \emph{arXiv preprint arXiv:2610.17907}, 2026.
\end{thebibliography}

\end{document}
